\chapter{PENDAHULUAN}
\label{chap:pendahuluan}

\paragraphtextstyle

\section{LATAR BELAKANG}
\label{sec:latarbelakang}

Ditulis latar belakang laporan magang. Latar belakang menjelaskan permasalahan, mengapa masalah tersebut penting untuk dipecahkan dan lain-lain.

\section{PERMASALAHAN}
\label{sec:permasalahan}

Permasalah yang diuraikan berupa permasalahan yang dihadapi selama proses magang di tempat Anda.

\section{TUJUAN}
\label{sec:tujuan}

Tujuan diuraikan dengan contoh sebagai berikut:
\begin{enumerate}[label=\alph*.]
  \item Memberikan gambaran dunia kerja yang sebenarnya kepada mahasiswa sebagai bekal untuk kemudian hari.
\end{enumerate}

\section{MANFAAT MAGANG}
\label{sec:manfaat}

\subsection{Bagi Mahasiswa}
\begin{enumerate}[label=\alph*.]
  \item Dapat memperoleh kesempatan untuk menerapkan ilmu pengetahuan yang telah diperoleh dalam perkuliahan untuk diterapkan dalam lapangan kerja.
  \item dan seterusnya
\end{enumerate}

\subsection{Bagi Perguruan Tinggi}
\begin{enumerate}[label=\alph*.]
  \item Mengenalkan dunia profesional terkait alat kesehatan kepada mahasiswa agar tercipta mahasiswa yang unggul, berpengalaman dan berkualitas sehingga siap terjun ke dunia kerja.
  \item dan seterusnya
\end{enumerate}

\subsection{Bagi Perusahaan}
\begin{enumerate}[label=\alph*.]
  \item Mahasiswa dapat membantu pekerjaan/proyek.
  \item dan seterusnya
\end{enumerate}
